\section{Proposed RC-QASyn Framework}
\label{sec:method}

\rcqasyn{} is an intended architecture for constructing a synthetic dataset
under explicit quality and risk requirements. Its core comprises scoped
record- and dataset-level evidence, calibration and fusion, and risk-aware
decision control around a replaceable generator. Localized repair,
privacy/fairness or domain-factual modules, and richer dataset optimization are
extensions that can be added when the core evidence justifies them. The
architecture is shown in \cref{fig:framework}; visual status distinguishes
implemented M1 infrastructure, M1.5 diagnostics, proposed core components, and
later extensions.

\begin{figure}[t]
  \centering
  \begin{tikzpicture}[
    node distance=3.4mm,
    every node/.style={font=\small, align=center, minimum height=6.5mm,
      text width=0.70\linewidth, rounded corners=2pt},
    implemented/.style={draw=black, fill=blue!8, thick},
    diagnostic/.style={draw=orange!75!black, fill=orange!10, thick},
    proposed/.style={draw=gray!75, fill=gray!6, dashed, thick},
    flow/.style={-{Latex[length=2mm]}, thick}
  ]
    \node[implemented] (spec) {Dataset specification and TRAIN-derived context};
    \node[implemented, below=of spec] (gen) {Versioned generation and persisted raw responses};
    \node[implemented, below=of gen] (val) {Record-level validators\\\footnotesize deterministic M1 baseline; broader core set proposed};
    \node[diagnostic, below=of val] (diag) {Dataset-level evidence\\\footnotesize copying and support measured in M1.5; not acceptance criteria};
    \node[implemented, below=of diag] (evidence) {Scoped heterogeneous evidence\\\footnotesize current contracts; explicit scope proposed};
    \node[proposed, below=of evidence] (fusion) {Calibration and evidence fusion};
    \node[proposed, below=of fusion] (control) {Risk-aware acceptance decision};
    \node[proposed, below=of control] (action) {Accept \;/\; regenerate \;/\; reject\\\footnotesize localized repair is a later extension};
    \node[proposed, below=of action] (dataset) {Richer dataset optimization\\\footnotesize later extension if justified};
    \node[implemented, below=of dataset] (audit) {Provenance, artifacts, and human-audit infrastructure};
    \node[implemented, below=of audit] (utility) {Intrinsic evaluation and downstream utility};

    \draw[flow] (spec) -- (gen);
    \draw[flow] (gen) -- (val);
    \draw[flow] (val) -- (diag);
    \draw[flow] (diag) -- (evidence);
    \draw[flow] (evidence) -- (fusion);
    \draw[flow] (fusion) -- (control);
    \draw[flow] (control) -- (action);
    \draw[flow] (action) -- (dataset);
    \draw[flow] (dataset) -- (audit);
    \draw[flow] (audit) -- (utility);
  \end{tikzpicture}
  \caption{Intended \rcqasyn{} pipeline. Solid blue nodes denote implemented
  M1 infrastructure, orange denotes M1.5 diagnostic analysis, and dashed gray
  nodes denote unevaluated proposed components. The central proposed path is
  scoped evidence, calibration/fusion, and risk-aware decision making;
  localized repair and richer optimization are extensions.}
  \label{fig:framework}
\end{figure}


\subsection{Specification and Generation}

The specification makes the schema, target, categorical domains,
missing-value policy, class condition, generation objective, and provenance
requirements explicit. Generation uses only a TRAIN-derived profile and,
when enabled, deterministic class-aware TRAIN examples. Prompts and provider
settings are versioned. Raw responses are preserved before parsing, so a
transport or parse failure remains observable instead of being silently
repaired. This stage and its reproducibility machinery are implemented.

The generator is not assumed to certify its own output. A prompt may describe
observed ranges or request target consistency, but compliance is established
only by evidence produced after generation. This separation prevents prompt
instructions from being mistaken for enforced constraints.

\subsection{Independent Validators and Heterogeneous Evidence}

Each validator should answer a narrow question and emit a structured evidence
item rather than a final global judgment. Deterministic schema, domain,
missing-value, target-condition, observed-range, and duplicate checks exist in
the current baseline pipeline. The core proposed evaluator set adds
independent statistical-support and task-relevant logical or semantic
evidence. Factual verification, privacy, and fairness require domain- and
use-specific definitions and are treated as optional modules, not as
implemented or guaranteed capabilities in PAPER v0.1.

Validator independence is an experimental objective, not an assumption that
errors are statistically independent. Two evaluators may share training data,
model families, prompts, or representations and therefore make correlated
errors. Provenance and evaluator identity are retained so calibration and
ablation can measure such dependence rather than conceal it.

\subsection{Record-Level and Dataset-Level Evidence}
\label{sec:record-dataset-evidence}

Record-level evidence includes schema correctness, required-field presence,
categorical validity, target consistency, local statistical support, and
logical or semantic plausibility. It supports statements about one candidate
or one of its fields. Dataset-level evidence instead includes exact and near
redundancy, prompt-reference copying, support and class-conditional coverage,
distributional fidelity, subgroup representation, and downstream utility.

The distinction changes how evidence is acted upon. A target mismatch is a
candidate defect; adding many copies of a correct minority example is a
collection-level failure. A record can therefore satisfy record-level validity
while contributing little new support to the dataset. \rcqasyn{} proposes to
retain both views and to prevent a dataset-level redundancy finding from being
misreported as record invalidity.

\subsection{Calibration and Evidence Fusion}

The proposed calibration layer maps evaluator outputs to empirically supported
reliability statements on held-out, human-labelled development data. Fusion
then combines evidence without erasing scope, missingness, or provenance. A
field-level hard violation may dominate an acceptance decision, while
uncertain semantic evidence may call for review or regeneration. Dataset-level
coverage evidence should influence selection among otherwise acceptable
records rather than retroactively redefine their structural validity.

This draft deliberately does not specify a universal weighted sum. Calibration
targets, dependence handling, fusion rules, and risk guarantees must be frozen
and evaluated after the human audit. The current B3 score is uncalibrated and
is not an implementation of this layer.

\subsection{Risk-Aware Decision Control}

Given calibrated evidence and declared risk constraints, the intended
controller selects whether a candidate should be accepted, regenerated, or
rejected under a frozen loss and operating policy. The action contract also
reserves repair, but the core empirical question is whether calibrated scoped
evidence yields more reliable acceptance decisions than the paired baselines.
No action-selection policy or calibrated risk guarantee is implemented in the
present milestone.

\subsection{Dataset Evidence and Audit}

Dataset-level evidence evaluates marginal redundancy, coverage, fidelity, and
prospective utility relative to $D_{\mathrm{syn}}^{(t)}$. Keeping this evidence
separate from record validity is part of the core formulation. A future
controller may use it to defer or decline a valid but redundant record, but
thresholds and objectives must be derived without access to REAL TEST labels
or statistics.

M1.5 implements diagnostic measurements of exact copying and TRAIN-support
distance, but these measurements do not yet control acceptance. The audit
layer is already partly implemented: run manifests, resolved configuration,
prompt and seed plans, provider traces, candidates, decisions, evidence,
metrics, copying/support diagnostics, and blinded human-review identities are
persisted for analysis.

\subsection{Later-Stage Extensions}

Localized repair could be useful when a defect and its dependency set can be
identified and revalidated; complete regeneration remains the comparison when
localization is unsafe or uneconomical. Richer dataset optimization could
select among valid candidates using multiple marginal-contribution objectives.
Neither mechanism is necessary to test the core record/dataset evidence and
calibrated-decision hypothesis. No repair operator or dataset optimizer is
currently implemented, and neither is presented as an established advantage.

\subsection{Implementation Status Summary}

\begin{table}[t]
  \centering
  \caption{Framework status in PAPER v0.1. ``Development'' denotes pilot
  diagnostics, not final benchmark validation.}
  \label{tab:framework-status}
  \begin{tabularx}{\linewidth}{@{}lXl@{}}
    \toprule
    Component & Current repository support & Research status \\
    \midrule
    Specification and generation & Versioned configs/prompts, TRAIN-only
      context, persisted raw responses and provenance & Core---implemented \\
    Baseline validators & Deterministic B2 checks and one B3 critic over a
      shared candidate pool & Core---implemented \\
    Copying/support analysis & Exact-copy categories and TRAIN-distance
      summaries without acceptance thresholds & Core---development \\
    Human audit & Blinded export/import, overlap, agreement, and adjudication
      infrastructure; labels incomplete & Core---results pending \\
    Explicit evidence scope & FIELD/RECORD/SUBGROUP/DATASET interface & Core---proposed \\
    Calibration and fusion & Reliability estimation and heterogeneous fusion & Core---proposed \\
    Risk-aware controller & Evidence-based acceptance policy & Core---proposed \\
    Localized repair & Dependency-aware edit and revalidation & Extension \\
    Richer dataset optimization & Multi-objective marginal-contribution
      selection & Extension \\
    \bottomrule
  \end{tabularx}
\end{table}
